\chapter{Networking: nornNET, nornWIFI and nornHTTP}
\label{ch:net}

The network modules are separate repositories with their own READMEs and plans; this
chapter is what an application writer needs to use them from nornRTOS. The design
sketch, protocol scope and test record are in \code{nornNET/doc/plan.md},
\code{nornWIFI/doc/plan.md} and \code{nornHTTP/docs/} (contracts, measurements, plan).

\section{nornNET}

nornNET is a small TCP/IP stack: Ethernet, ARP, IPv4, ICMP, UDP, TCP (RFC 9293 with slow
start, fast retransmit, delayed ACK, Nagle, zero window probing and out of order
storage), a DHCP client and a DNS resolver, loosely modelled on microps~\cite{microps} in
structure. Packet buffers are fixed size blocks from a pool, owned by whoever holds the
pointer; there is one net task, one re-entrant lock, and no heap in the data path. It is
verified against the Linux kernel as the peer, over a TAP device, QEMU's user mode
network and a serial line (SLIP) on the ESP32 board, and it runs on the Mega 2560 over
SLIP at the line rate.

\subsection{Attaching it to an application}

In \code{app.mk}:
\begin{lstlisting}[language={}]
NORNNET_DIR = /path/to/nornNET
MODULES    += $(NORNNET_DIR)/nornnet.mk
\end{lstlisting}
and an \code{nn\_config.h} next to the application with at least
\code{NN\_CFG\_TASK\_PRIO}, the net task's priority, above every task that uses the
stack. Every \code{NN\_CFG\_*} option and its default is in \code{nn/nn\_cfg.h}. In
\code{main}, after \code{rtos\_init}: \code{nn\_init()}, register the devices, add the
interfaces with \code{nn\_ifaceAdd} (with an address, or empty and \code{nn\_dhcpStart}
from a task), then \code{rtos\_start()}.

\subsection{Sockets}

Sockets are caller owned objects; for TCP the caller also owns the send and receive
rings. Each call is either blocking with a timeout, from a task, or a callback that runs
in the net task. The lock is re-entrant for the net task, so a callback may call the
socket API.

\begin{lstlisting}
nn_err_t nn_udpOpen(nn_udp_t *pSock);
nn_err_t nn_udpBind(nn_udp_t *pSock, nn_ip4_t addr, uint16_t port);
nn_err_t nn_udpConnect(nn_udp_t *pSock, nn_ip4_t addr, uint16_t port);
nn_err_t nn_udpSendTo(nn_udp_t *pSock, const void *p, size_t n, nn_ip4_t to, uint16_t port);
nn_err_t nn_udpSend(nn_udp_t *pSock, const void *p, size_t n);
int32_t  nn_udpRecvFrom(nn_udp_t *pSock, void *p, size_t n, nn_ip4_t *pFrom,
                        uint16_t *pPort, nn_ms_t timeout);
void     nn_udpOnReadable(nn_udp_t *pSock, nn_udpReadableFunc_t fn, void *arg);
nn_err_t nn_udpClose(nn_udp_t *pSock);

nn_err_t nn_tcpOpen(nn_tcp_t *pSock, void *pRxRing, uint16_t rxSize, void *pTxRing, uint16_t txSize);
nn_err_t nn_tcpBind(nn_tcp_t *pSock, nn_ip4_t addr, uint16_t port);
nn_err_t nn_tcpListen(nn_tcp_t *pSock);
nn_err_t nn_tcpAccept(nn_tcp_t *pListen, nn_tcp_t *pNew, nn_ms_t timeout);
nn_err_t nn_tcpConnect(nn_tcp_t *pSock, nn_ip4_t addr, uint16_t port, nn_ms_t timeout);
int32_t  nn_tcpSend(nn_tcp_t *pSock, const void *p, size_t n, nn_ms_t timeout);
int32_t  nn_tcpRecv(nn_tcp_t *pSock, void *p, size_t n, nn_ms_t timeout);
nn_err_t nn_tcpClose(nn_tcp_t *pSock);
nn_err_t nn_tcpAbort(nn_tcp_t *pSock);
void     nn_tcpOnEvent(nn_tcp_t *pSock, nn_tcpEventFunc_t fn, void *arg);

nn_err_t nn_dhcpStart(nn_dhcp_t *pD, nn_iface_t *pIf, nn_dhcpChangeFunc_t fn, void *arg);
nn_err_t nn_dnsResolve(const char *name, nn_ip4_t *pAddr, nn_ms_t timeout);
\end{lstlisting}

Errors are negative \code{nn\_err\_t} values: \code{NN\_ERR\_NOMEM} (no buffer),
\code{NN\_ERR\_TIMEOUT}, \code{NN\_ERR\_DOWN} (the device or the address is gone),
\code{NN\_ERR\_RESET}, \code{NN\_ERR\_REFUSED}, \code{NN\_ERR\_CLOSED},
\code{NN\_ERR\_PENDING} (a non-blocking call accepted) among them. Timeouts in nornNET
are milliseconds (\code{nn\_ms\_t}).

\subsection{Links that come and go}

A driver whose carrier can disappear (WiFi, a cable) reports it with
\code{nn\_netdevLinkChanged(dev, up)} from any context. On a loss the net task drops
the device's ARP entries; on a return every DHCP client on the device confirms its lease
at once, keeping the address meanwhile, so an outage on the same network does not
disturb TCP connections. When an interface's address changes or goes, the TCP connections
bound to the old one end with \code{NN\_ERR\_DOWN}.

\section{nornWIFI}

Espressif ships the ESP32 radio as binary libraries written against a table of operating
system functions. nornWIFI implements that table over the nornCore primitives, adds the
chip glue (power domain, clocks, PHY calibration, interrupts), the WPA supplicant for
WPA2 personal, and a nornNET device. The result is a station that scans, joins, gets an
address by DHCP and carries TCP at 400 to 650\,kB/s each way, with nornCore as the only
kernel in the image.

\subsection{Using it}

\code{app.mk} adds both modules; \code{rtos\_config.h} needs \code{RTOS\_CFG\_HEAP 1},
\code{RTOS\_CFG\_EVT\_BITS 32}, \code{RTOS\_CFG\_EVGROUPS 1} and two spare task ids at
the top (\code{NW\_CFG\_TASK\_WIFI}, \code{NW\_CFG\_TASK\_SERVICE}; by default the two
highest). Then:

\begin{lstlisting}
nw_init(NW_HEAP_SPARE_START, NW_HEAP_SPARE_SIZE);    /* in main, before rtos_start */
...
nw_start();                                            /* in a task */
nn_ifaceAdd(nw_netdev(), &iface, zero, zero, zero);
nw_connect("ssid", "password", 20000);                 /* or nw_connectStart + events */
nn_dhcpStart(&dhcp, &iface, onDhcp, NULL);
\end{lstlisting}

From then on the station keeps the link up by itself (\code{NW\_CFG\_RECONNECT}: straight
back after a loss, then a delay doubling from 1 to 60 seconds) until
\code{nw\_disconnect}, and nornNET follows the link as described above.
\code{nw\_linkState}, \code{nw\_rssi}, \code{nw\_isConnected} and the events delivered
through \code{nw\_onEvent} (run in the service task, which must not call nornNET) tell
what happens. Measured on the board: three forced drops recover in about 160\,ms, with
DHCP confirmed within 100\,ms and a TCP connection intact.

\begin{lstlisting}
void      nw_init(void *heap, size_t heapSize);
esp_err_t nw_start(void);
esp_err_t nw_scan(wifi_ap_record_t *pRecords, uint16_t *pCount, uint32_t timeoutMs);
esp_err_t nw_connect(const char *ssid, const char *password, uint32_t timeoutMs);
esp_err_t nw_connectStart(const char *ssid, const char *password);
esp_err_t nw_disconnect(void);
bool      nw_isConnected(void);
uint8_t   nw_linkState(void);
int8_t    nw_rssi(void);
struct nn_netdev *nw_netdev(void);
void      nw_onEvent(nw_eventFunc_t cb, void *arg);
void      nw_getStats(nw_stats_t *pStats);
\end{lstlisting}

The radio libraries are Espressif's and keep their Apache~2.0 licence; the supplicant is
the upstream wpa\_supplicant as ported by Espressif (BSD). Both are listed in the
repository's \code{THIRD\_PARTY\_LICENSES.md}.

\section{nornHTTP}
\label{sec:nornhttp}

nornHTTP is an HTTP/1.1 client in portable C99 with nothing underneath it: no heap, no
global state, no recursion, no dependency on nornCore or nornNET. The application owns
every buffer and every context, so they can be static, and the library talks to the
network through three functions it is handed (send, receive, a millisecond clock). An
adapter for a connected nornNET TCP socket is part of the module; the same code runs
on Linux against the module's test server, which is where most of its verification
happens. It does GET, POST, PUT, DELETE, PATCH, HEAD and OPTIONS, custom headers,
request bodies from a buffer or streamed from a callback with a known length,
Content-Length, chunked and close-delimited responses, trailers, informational
responses, closed byte ranges with 206 checking, HTTP/1.0 servers, and keep-alive. It
does not do TLS, redirects, compression, cookies, DNS or HTTP/2; TLS is a transport the
application would supply.

The response body is streamed: the parser hands each span of decoded bytes to a
callback as it arrives, so a download of megabytes needs only the receive buffer the
application chose (the ESP32 example below uses 512 bytes). The size of a response does
not decide the RAM cost; the longest header line and the longest request head do.

\subsection{Using it}

\code{app.mk} adds the sources and the adapter; the \code{wifi\_http} example in
nornRTOS has the lines. With a socket already connected (\code{nn\_tcpConnect}), one
request is:

\begin{lstlisting}
static char line[256], request[256];          /* longest header line, request head */
static uint8_t receive[512], body[128];       /* one recv at a time, a small body */
static norn_http http;

norn_http_net_context net = { &socket, now_ms, NULL };   /* now_ms wraps nn_os_now */
norn_http_transport transport;
norn_http_config config = {0};
norn_http_body_buffer sink = { body, sizeof(body), 0 };
norn_http_request req;
norn_http_response resp;

norn_http_net_transport_init(&transport, &net);
norn_http_parser_config_init(&config.parser, line, sizeof(line));
config.parser.on_body = norn_http_buffer_body;    /* or your own callback */
config.parser.callback_context = &sink;
config.receive_buffer = receive;   config.receive_capacity = sizeof(receive);
config.request_buffer = request;   config.request_capacity = sizeof(request);
config.timeout_ms = 5000;                         /* the whole request and response */
norn_http_init(&http, &config, &transport);

norn_http_request_init(&req, "GET", "/hello", "192.168.1.10:18080", NULL, 0);
if(norn_http_execute(&http, &req, &resp) == NORN_HTTP_OK)
    ...                        /* resp.status, sink.length bytes in body, resp.reusable */
\end{lstlisting}

A POST adds \code{norn\_http\_add\_header(\&req, "Content-Type", "application/json")}
and \code{norn\_http\_set\_body(\&req, json, length)}, or
\code{norn\_http\_set\_body\_source} with a callback that fills the request buffer span
by span when the body is not in RAM as a whole. A large download replaces the body
callback with one that writes each span to its destination, or checks it, and stores
nothing. \code{resp.reusable} says whether the same socket may carry the next request;
\code{req.close} asks the server to close after this one.

Every function returns a \code{norn\_http\_result}; zero is success, the HTTP status is
a separate matter in \code{resp.status}. After a negative result from
\code{norn\_http\_execute} the connection is in an unknown state: abort the socket and
call \code{norn\_http\_init} again (the context answers \code{NORN\_HTTP\_STALE} until
then). An argument error (\code{INVALID}, or \code{LIMIT} from the request builder) is
reported before anything is sent and leaves the connection as it was. The full
contract, including what the parser rejects and why, is \code{nornHTTP/docs/contracts.md};
the API reference is generated from \code{norn\_http.h} together with the kernel's.

\subsection{What it costs}

Measured on the NodeMCU-32S with the \code{wifi\_http} example (GET, JSON POST, a
64\,KiB download through the 512-byte buffer with every byte checked, keep-alive, then
Connection: close): the whole image grows by 7.0\,KB of flash, of which about 5.3\,KB is
the library and 0.2\,KB the nornNET adapter. The library has no static RAM; the
context is 184 bytes, a request 64 and a response 56, and the example's buffers add
1.2\,KB. The main task's stack high-water went from 1.3\,KB after the Wi-Fi join and
TCP connect to at most 1.5\,KB after the four requests, of 4\,KB; the library's own
frames are small (the deepest chain inside it is about 240 bytes), the rest is the TCP
path and the kernel running on the caller's stack. Details and the method are in
\code{nornHTTP/docs/measurements.md}.
